\documentclass[lettersize,journal]{IEEEtran}

\usepackage{amsmath,amsfonts}
\usepackage{graphicx}
\usepackage{cite}
\usepackage{algorithm}
\usepackage{algorithmic}
\usepackage{array}
\usepackage{url}
\usepackage{textcomp}
\usepackage{stfloats}
\usepackage{placeins}

\begin{document}


% ===== Title / Authors (replace your current \title and \author block) =====
\title{GIS-Constrained Wake-Aware Wind Farm Layout Optimization and Scalable Expansion: A Case Study and National Demonstration}

\author{Soheil Ayati and Marc Steffgen%
\thanks{Soheil Ayati and Marc Steffgen are with TH Köln (University of Applied Sciences),  Faculty of Computer Science and Engineering Science, Campus Gummersbach, Germany.}%
\thanks{Sustainable Energy Systems, Project Paper, submitted: Jan.~2026.}%
}

\maketitle

\begin{abstract}
Wake effects can significantly reduce wind farm energy yield, while land-use and planning constraints limit where turbines can be placed. This paper presents a proof-of-concept workflow for wake-aware layout optimization using the Hill of Towie wind farm as a case study. Module~1 builds a SCADA-based baseline and quantifies wake losses with a Jensen/Park model. Module~2 optimizes turbine positions under Geographic Information System (GIS) and spacing constraints and applies a controlled, sequential expansion strategy. Module~3 validates performance comparability and feasibility. In addition, a coarse ``Project Deutschland 2030'' demonstration shows how the same logic can be transferred to a national scale using clustering and proxy models.
\end{abstract}


\begin{IEEEkeywords}
Wind farm layout optimization, wake model, Jensen/Park, GIS constraints, Differential Evolution.
\end{IEEEkeywords}


% ============================================================
% Problemstatement
% ============================================================

\section{Introduction}

Wind energy is expected to play a central role in Germany’s future energy system. In 2024, wind power accounted for approximately 30\,\% of national electricity generation, with an installed on shore capacity of about 59~GW as of 2023. To meet climate and energy security targets, German onshore wind capacity is planned to increase to 115~GW by 2030 and further to approximately 160~GW by 2040 \cite{bundestag2023wind}.

\begin{figure}[h]
\centering
\includegraphics[width=\columnwidth]{German Targets Windturbines.png}
\caption{Historical and projected development of onshore wind energy capacity in Germany. Gray bars represent historical annual capacity additions, while white bars indicate expected future additions based on EEG auction volumes. The blue shaded area shows the cumulative installed capacity. The solid line indicates the legally defined expansion targets according to the German Renewable Energy Act (EEG 2021/2023), illustrating the required acceleration of wind energy deployment toward 2030. \cite{bundestag2023wind}}
\label{fig:wake_layout}
\end{figure}


This expansion is increasingly shaped by regulatory frameworks and, in particular, by societal acceptance of new wind energy projects. While several studies indicate that Germany possesses sufficient technical land potential for large-scale wind energy deployment, the practical realization of new wind farms is often constrained by public acceptance and planning processes. As a result, there is a strong incentive for operators and planners to maximize energy production within already designated sites. Improving the efficiency of existing wind farm layouts reduces the need for additional turbines and new park developments, thereby mitigating acceptance-related conflicts while simultaneously improving economic performance. Consequently, wake-aware optimization becomes a key lever to increase energy yield without expanding the spatial footprint of wind energy installations \cite{acceptanceOnshoreWind2021}.

In such layouts, aerodynamic wake effects become a dominant limiting factor. Turbines operating in the wake of upstream machines experience reduced wind speeds and increased turbulence, resulting in significant energy losses and elevated structural loads. Reported wake-induced annual energy losses typically range between 10\,\% and 20\,\% for onshore wind farms, with increasing relevance at higher turbine densities \cite{archer2018wakeloss}.

Despite their importance, wake interactions are often only implicitly considered during wind farm planning through conservative spacing rules or simplified heuristics. These approaches do not systematically optimize turbine placement at the park level and can lead to avoidable energy losses and inefficient land use \cite{archer2018wakeloss}.

The problem addressed in this work is therefore the lack of a systematic, wake-aware layout optimization approach that maximizes energy yield within fixed land constraints. The objective is to improve wind farm performance by reducing wake losses and enabling more efficient use of limited available space.

% ============================================================
% State of the Art
% ============================================================

\section{State of the art}

Wind farm layout design has traditionally relied on rule-based turbine spacing and heuristic placement strategies, with wake effects often evaluated only indirectly using simplified engineering models. While such approaches are computationally efficient, they do not systematically optimize turbine placement at the park level and may lead to avoidable wake-induced energy losses. As wind farms increase in size and density, wake-aware wind farm layout optimization (WFLO) has therefore become an active research field\cite{dtuStateOfArtWFLO}.

State-of-the-art WFLO approaches typically combine analytical wake models with metaheuristic or gradient-based optimization methods to maximize annual energy production under spatial and regulatory constraints \cite{herbertacero2014review}\cite{thomas2023wflo}. Wake modeling approaches range from fast engineering models, such as the Jensen (Park) model, to more advanced Gaussian wake formulations and high-fidelity CFD/LES simulations. While higher-fidelity models provide improved physical accuracy, their computational cost often limits their applicability in iterative optimization \cite{bastankhah2014gaussian}. Consequently, simplified engineering wake models remain widely used for early-stage design and proof-of-concept optimization studies, where relative layout comparisons are of primary interest\cite{katic1986park}.

Within this context, the present work follows established practice by employing a computationally efficient engineering wake model to demonstrate the potential of systematic, park-wide layout optimization under fixed land constraints.


% ============================================================
% Conzept Overview
% ============================================================

\section{Concept Overview}

The proposed concept aims to improve wind farm energy yield by explicitly considering wake effects during turbine layout optimization. The approach is structured into three modules:
\begin{itemize}
    \item Module 1: Baseline modeling and wake loss quantification.
    \item Module 2: Wake-aware layout optimization under land constraints.
    \item Module 3: Validation of optimization results.
\end{itemize}

The concept is implemented and evaluated using the Hill of Towie onshore wind farm in Scotland as a representative case study. The site was selected due to its typical park size and the availability of extensive operational data. To ensure feasibility within the project scope, the wind farm is modeled in two dimensions using simplified physical assumptions. The concept is explicitly designed as a proof of concept, prioritizing transparency, robustness, and computational efficiency.


% ============================================================
% MODULE 1 
% ============================================================
\section{Module~1: Baseline Modeling and Results}

Module~1 establishes a physically consistent baseline model of the existing Hill of Towie wind farm and quantifies wake-induced energy losses under real operating conditions. SCADA data from the year 2023, including turbine-level wind speed, yaw position, and power output at a 10-minute resolution, is used as the reference.

Each turbine is modeled using a simplified power curve that converts effective hub-height wind speed into electrical power. Wake effects are represented using the Jensen (Park) wake model, which approximates wakes as linearly expanding velocity deficits downstream of each turbine. 
\begin{equation}
\delta_{u \rightarrow d}(x)
=
\left(
1 - \sqrt{1 - C_T}
\right)
\left(
\frac{D}{D + 2 k x}
\right)^2
\end{equation}

\begin{equation}
U_d(t)
=
U_{\infty}(t)
\left(
1 - \sqrt{
\sum_{u \in \mathcal{U}(d,t)}
\delta_{u \rightarrow d}^2
}
\right)
\end{equation}

To reduce model complexity, the wind farm is represented in a two-dimensional horizontal plane and partial wake overlaps are neglected. A global calibration factor is applied to account for unmodeled system losses and to align the modeled output with measured production.

For each time step, three quantities are computed: measured power output from SCADA data, modeled power output including wake effects,
\begin{equation}
P_{\text{wake}}(t)
=
\sum_{d=1}^{N}
P_{\text{pc}}\!\left(U_d(t)\right)
\end{equation}
 and a hypothetical no-wake power output assuming free-stream conditions at all turbines.
 \begin{equation}
P_{\text{pot}}(t)
=
N \, P_{\text{pc}}\!\left(U_{\infty}(t)\right)
\end{equation}
Annual energy production is obtained by integrating the time series over the full year.
\begin{equation}
E
=
\sum_{t}
P(t)\,\Delta t
\end{equation}


The results show good agreement between modeled and measured energy production, indicating that the simplified physical model captures the dominant wake-related effects. The analysis reveals wake-induced annual energy losses of approximately 9--14\,\%, while the no-wake scenario represents an upper bound on achievable production. The modeled wake-affected energy yield serves as the baseline reference for the optimization performed in Module~2.

\begin{figure}[H]
\centering
\includegraphics[width=\columnwidth]{Modul1_result_graph.png}
\caption{Comparison between the annual power output predicted by the physical wake model using wind rose data and the measured power output of the Hill of Towie wind farm for the same year.}
\label{fig:wake_layout}
\end{figure}


% ============================
% MODULE 2 (IEEE paper text)
% ============================

\section{Module 2: GIS-Constrained Layout Optimization and Dynamic Expansion}
\label{sec:module2}

Module~2 maximizes annual energy production (AEP) by (i) repositioning the existing $N_0=21$ turbines and (ii) adding new turbines under strict geospatial and engineering constraints. The final approach (\textit{Implementation~2}) combines a fast Jensen/Park wake model with a hybrid optimization pipeline and a sequential expansion strategy that stops automatically once additional turbines become inefficient or economically unattractive. An earlier prototype (\textit{Implementation~1}) used a restricted micrositing search and a fixed expansion of two turbines; it is reported only for comparison.

\subsection{Inputs and Coordinate System}
The optimizer uses turbine metadata (rotor diameter $D$, rated power $P_{\text{rated}}$, coordinates), a weighted wind rose with bins $\{(u_b,\theta_b,p_b)\}_{b=1}^{B}$, and geospatial constraints extracted from OpenStreetMap (OSM) features.

\subsection{Wake-Aware AEP Estimation}
To enable thousands of layout evaluations, wake interactions are approximated using the Jensen/Park model \cite{shakoor2016jensenreview,katic1987jensen}. For each wind direction bin, turbines are sorted in the downwind direction and downstream turbines receive a reduced effective wind speed due to upstream wakes. Power is then computed using a standard cubic-region power curve with cut-in, rated, and cut-out speeds (implemented in code as a piecewise model; parameters listed in Table~\ref{tab:m2_params}).

The wind-rose AEP estimator is
\begin{equation}
\widehat{\mathrm{AEP}}_{\text{raw}}(\mathbf{x}) = 
\sum_{b=1}^{B} \left[
p_b \cdot 8760 \cdot \sum_{j=1}^{N} \widehat{P}\!\left(u_{j,b}^{\text{eff}}(\mathbf{x})\right)
\right],
\label{eq:m2_aep_raw}
\end{equation}
where $u_{j,b}^{\text{eff}}$ is the wake-reduced wind speed at turbine $j$ under bin $b$, and $\widehat{P}(\cdot)$ denotes the implemented power-curve model.

\subsection{Calibration to the SCADA Baseline (Module 1)}
Because \eqref{eq:m2_aep_raw} aggregates a wind rose instead of replaying the full SCADA time series, its absolute magnitude is calibrated to the validated Module~1 baseline AEP. Using the original (as-built) layout $\mathbf{x}_0$:
\begin{equation}
\eta = \frac{\mathrm{AEP}_{\text{M1}}}{\widehat{\mathrm{AEP}}_{\text{raw}}(\mathbf{x}_0)},
\qquad
\mathrm{AEP}(\mathbf{x}) = \eta\cdot \widehat{\mathrm{AEP}}_{\text{raw}}(\mathbf{x}).
\label{eq:m2_calibration}
\end{equation}
This ensures all optimized layouts are reported on the same absolute energy scale as the SCADA-validated baseline.

Wake loss is reported relative to the Module~1 no-wake potential. For expanded layouts $(N>N_0)$, the potential is scaled linearly to isolate wake-interaction efficiency:
\begin{equation}
\mathrm{AEP}_{\text{pot}}(N) \approx \mathrm{AEP}_{\text{pot}}(N_0)\cdot\frac{N}{N_0}.
\label{eq:m2_potential_scale}
\end{equation}

\subsection{GIS and Engineering Constraints}
Constructibility is enforced using (i) forbidden regions created by buffering OSM road and building layers, (ii) a farm boundary polygon with an additional legal buffer, (iii) minimum inter-turbine spacing, and (iv) an \emph{infill zone} that prevents unrealistic edge scattering by forcing new turbines to be placed inside a contracted allowed region. Rather than repeating all constraint equations, the final parameter set is summarized in Table~\ref{tab:m2_params}. Infeasible candidates are rejected or assigned large penalties so they are never selected by the optimizer.

\begin{table}[H]
\caption{Key optimization, wake, and constraint parameters used in Module~2.}
\label{tab:m2_params}
\centering
\renewcommand{\arraystretch}{1.1}
\begin{tabular}{lc}
\hline
\textbf{Parameter} & \textbf{Value} \\
\hline
Wake model & Jensen/Park \cite{shakoor2016jensenreview,katic1987jensen} \\
Wake decay $k$ & 0.075 \\
Thrust coefficient $C_T$ & 0.78 \\
Power-curve scaling $\gamma$ & 0.45 \\
Cut-in / rated / cut-out & 3 / 12 / 25 m/s \\
Road buffer $R_{\text{road}}$ & 120 m \\
Building buffer $R_{\text{bldg}}$ & 200 m \\
Boundary buffer (legal) & 300 m \\
Infill contraction $R_{\text{edge}}$ & 50 m \\
Min spacing (existing--existing) & $3.5D$ \\
Min spacing (any pair with new) & $5.0D$ \\
Stop: marginal gain & 3.4 GWh \\
Stop: wake-loss increase & $+1.0$ pp \\
\hline
\end{tabular}
\end{table}

\subsection{Implementation 1 (Prototype: Micrositing + Fixed Expansion)}
Implementation~1 served as an intermediate baseline to verify that the wake-aware objective and the GIS penalty system improve AEP under realistic movement limits. In Step~1, turbines were \emph{microsited} by restricting each turbine to a small search box around its as-built position (order of a few hundred meters), which represents typical planning freedom when repowering or refining a permitted site. The optimizer then improved the layout using the same wake-based AEP estimator, but only within this constrained neighborhood.

In Step~2, two additional turbines were placed \emph{with a fixed turbine count} rather than an adaptive stopping rule. Candidate locations were searched inside the allowed region while enforcing spacing and OSM-based exclusion zones, and the best feasible additions were accepted. This prototype demonstrates the benefit of wake-aware optimization without changing the farm capacity too aggressively; its results are reported in Table~\ref{tab:m2_results} for comparison.

\subsection{Optimization Workflow (Implementation 2)}
\subsubsection{Step 1: Repositioning the existing turbines}
The existing $N_0$ turbines are optimized inside the legal boundary using a three-stage hybrid pipeline:
\begin{enumerate}
\item \textbf{Global search (DE):} Differential Evolution explores the non-convex wake-loss surface.
\item \textbf{Micro-polish:} small discrete moves accelerate improvement near good basins.
\item \textbf{Gradient polish (L-BFGS-B):} refines the best candidate under box/boundary constraints.
\end{enumerate}
This combination yields robust global exploration while retaining fast local convergence.

\subsubsection{Step 2: Dynamic marginal expansion (sequential turbine addition)}
New turbines are added one at a time. At each iteration, the position of the next turbine is optimized while existing turbines remain fixed. The candidate is accepted only if it satisfies all GIS and spacing constraints (including the infill rule). The expansion stops automatically if either (i) the marginal energy gain drops below a fixed cutoff or (ii) wake loss rises more than a fixed tolerance above the baseline.

\begin{algorithm}[H]
\caption{Dynamic marginal expansion (Module~2, Implementation~2)}
\label{alg:m2_expansion}
\begin{algorithmic}[1]
\STATE \textbf{Input:} optimized layout $\mathbf{x}$, calibration factor $\eta$, baseline wake loss $\mathcal{L}_0$
\STATE $\mathrm{AEP}_0 \leftarrow \mathrm{AEP}(\mathbf{x})$
\FOR{$k=1,\dots,K_{\max}$}
  \STATE Optimize candidate position $\mathbf{p}_k$ using DE subject to GIS + spacing constraints
  \IF{$\mathbf{p}_k$ violates GIS/spacing/infill constraints}
    \STATE \textbf{break}
  \ENDIF
  \STATE $\mathbf{x}' \leftarrow [\mathbf{x}; \mathbf{p}_k]$ \hfill (tentative layout)
  \STATE $\mathrm{AEP}_k \leftarrow \mathrm{AEP}(\mathbf{x}')$
  \STATE $\Delta\mathrm{AEP}_k \leftarrow \mathrm{AEP}_k - \mathrm{AEP}_{k-1}$
  \STATE Compute wake loss $\mathcal{L}_k$ using $\mathrm{AEP}_{\text{pot}}(N)$ from \eqref{eq:m2_potential_scale}
  \IF{$\Delta\mathrm{AEP}_k < 3.4~\mathrm{GWh}$ \textbf{or} $\mathcal{L}_k > \mathcal{L}_0 + 1.0~\text{pp}$}
    \STATE \textbf{break}
  \ENDIF
  \STATE $\mathbf{x} \leftarrow \mathbf{x}'$ \hfill (accept turbine)
\ENDFOR
\STATE \textbf{Output:} final expanded layout $\mathbf{x}$
\end{algorithmic}
\end{algorithm}

\subsection{Results and Discussion}
Table~\ref{tab:m2_results} summarizes the outcomes on the calibrated AEP scale. The final approach (Implementation~2) achieves a substantially higher AEP by dynamically identifying the site capacity while maintaining wake loss close to the baseline due to the explicit guard rails and the infill constraint. Fig.~\ref{fig:m2_finalmap} visualizes the resulting layout for Implementation~2, including (i) repositioned original turbines, (ii) newly added turbines, and (iii) forbidden zones and boundary constraints.

\begin{figure}[H]
\centering
\includegraphics[width=\columnwidth]{f_final_map_module2.png}
\caption{Module~2 (Implementation~2) final layout map: optimized positions of the original 21 turbines, nine added turbines from dynamic expansion, and the GIS constraints (boundary and exclusion zones).}
\label{fig:m2_finalmap}
\end{figure}

\begin{table}[H]
\caption{Module~2 performance summary (calibrated to the Module~1 baseline).}
\label{tab:m2_results}
\centering
\renewcommand{\arraystretch}{1.1}
\begin{tabular}{lccc}
\hline
\textbf{Scenario} & $\mathbf{N}$ & \textbf{AEP (GWh)} & \textbf{Wake Loss (\%)} \\
\hline
Baseline (Module~1 reference) & 21 & 89.04 & 8.67 \\
Impl.~1: Reposition (micrositing) & 21 & 92.57 & 5.05 \\
Impl.~1: +2 turbines (fixed) & 23 & 100.67 & 5.72 \\
Impl.~2: Reposition (revamped) & 21 & 90.87 & 6.79 \\
Impl.~2: +9 turbines (dynamic) & 30 & 127.35 & 8.56 \\
\hline
\end{tabular}
\end{table}


% ============================================================
% MODULE 3 — Validation Framework (IEEE-style, non-duplicative)
% ============================================================

\section{Module 3: Validation Framework}
\label{sec:validation}

Module~3 verifies that the optimized layouts produced in Module~2 are (i) \emph{performance-credible} relative to the SCADA-derived baseline from Module~1, (ii) \emph{constructible} under GIS and spacing constraints, and (iii) \emph{methodologically justified} according to established Wind Farm Layout Optimization (WFLO) literature. To avoid repetition, this section focuses on the \emph{validation criteria and checks}, while the AEP estimator, calibration, and optimization workflow remain defined in Module~2.

% ------------------------------------------------------------
% Result Validation
% ------------------------------------------------------------

\subsection{Result Validation: Performance Credibility}
\label{subsec:result_validation}

\subsubsection{Baseline consistency check}
Module~2 evaluates layouts with a wind-rose AEP estimator and calibrates it to the SCADA-validated Module~1 baseline using the scalar factor in \eqref{eq:m2_calibration}. The primary credibility check is that the calibrated estimator reproduces the as-built baseline when evaluated at the original layout $\mathbf{x}_0$:
\begin{equation}
\varepsilon_{\mathrm{base}} \;=\;
\left|
\frac{\mathrm{AEP}(\mathbf{x}_0)-\mathrm{AEP}_{\mathrm{M1}}}{\mathrm{AEP}_{\mathrm{M1}}}
\right|.
\label{eq:m3_baseline_residual}
\end{equation}
In this study, $\varepsilon_{\mathrm{base}}$ is enforced to remain small (target $<0.5\%$), so that reported improvements are attributable to layout changes rather than an AEP scaling artifact.

\subsubsection{Validated improvement reporting}
After baseline alignment, layout improvements are interpreted via the calibrated results reported in Table~\ref{tab:m2_results}. We summarize performance using:
\begin{equation}
\Delta \mathrm{AEP} \;=\; \mathrm{AEP}(\mathbf{x}) - \mathrm{AEP}(\mathbf{x}_0),
\qquad
\Delta L_{\mathrm{pp}} \;=\; L(\mathbf{x}_0)-L(\mathbf{x}),
\label{eq:m3_gains}
\end{equation}
where $L(\cdot)$ is computed consistently with Module~2. For expanded layouts $(N>N_0)$, the no-wake potential scaling from \eqref{eq:m2_potential_scale} is applied to keep wake-loss values comparable across different turbine counts.

% ------------------------------------------------------------
% Constraint Validation
% ------------------------------------------------------------

\subsection{Constraint Validation: Constructibility and Compliance}
\label{subsec:constraint_validation}

\subsubsection{GIS feasibility and exclusion zones}
Constructibility is validated by confirming that all final turbine locations lie inside the allowed farm region and outside buffered exclusion zones derived from OpenStreetMap. Formally, with allowed region $\Omega$ and forbidden region $\mathcal{F}$ (roads/buildings buffers), the feasible set is
\begin{equation}
\mathcal{X}_{\mathrm{feas}} \;=\; \Omega \setminus \mathcal{F}.
\label{eq:m3_feasible_set}
\end{equation}
The final layouts are declared feasible if $\mathbf{p}_i \in \mathcal{X}_{\mathrm{feas}}$ for all turbines. In Implementation~2, an additional cohesion constraint is verified by checking that all newly added turbines satisfy the infill condition (new sites remain inside the contracted allowed region).

\subsubsection{Minimum spacing verification}
For all turbine pairs $(i,j)$, spacing compliance is validated with:
\begin{equation}
d_{ij} \;=\; \|\mathbf{p}_i-\mathbf{p}_j\|_2 \;\ge\; d_{\min},
\label{eq:m3_spacing}
\end{equation}
where $d_{\min}=3.5D$ is enforced for existing--existing pairs and $d_{\min}=5.0D$ for any pair involving a newly added turbine (Module~2 parameterization). These values are consistent with commonly cited practical ranges for onshore wind farm spacing \cite{luetkehus2013spacing}.

\subsubsection{Residential buffers and safety margin}
Residential setback conservatism is defended by comparing the applied building buffer to a simple topple-distance proxy:
\begin{equation}
d_{\mathrm{topple}} \approx 1.1H,
\label{eq:m3_topple}
\end{equation}
where $H$ is tip height. With representative geometry, $d_{\mathrm{topple}}\approx 165\,\mathrm{m}$, while the applied effective buffer ($\approx 350\,\mathrm{m}$) provides a conservative margin and aligns with setback discussions in planning literature \cite{rankl2024planning}.

% ------------------------------------------------------------
% Methodology Validation
% ------------------------------------------------------------

\subsection{Methodology Validation}
\label{subsec:methodology_validation}

\subsubsection{Global optimization choice (Differential Evolution)}
WFLO is non-convex with many local minima due to nonlinear wake interactions and discontinuous feasibility constraints; population-based metaheuristics are therefore widely used. Differential Evolution (DE) is employed as the global search engine, complemented by local refinement for efficient convergence \cite{herbertacero2014review}.

\subsubsection{Wake model choice (Jensen/Park)}
High-fidelity CFD/LES is computationally prohibitive inside iterative optimization. The Jensen/Park wake model is a common engineering workhorse in WFLO due to its analytical simplicity and speed, enabling thousands of objective evaluations \cite{shakoor2016jensenreview,katic1987jensen}.

\subsubsection{Sequential expansion as a planning proxy}
Implementation~2 expands the farm sequentially and stops when marginal benefit becomes too small or wake-loss degradation exceeds a tolerance. Such greedy/sequential strategies are frequently used as practical approximations to incremental investment decisions and can achieve competitive solutions at lower computational cost \cite{chen2016greedy}.

% ============================
% Error Analysis
% ============================

\section{Error Analysis and Uncertainty Assessment}
\label{sec:error_analysis}

This project is a proof-of-concept study; therefore, the reported AEP gains should be interpreted as \emph{relative improvements within a consistent modeling framework}, rather than exact predictions of absolute wind farm performance. This section documents the main uncertainty sources and provides transparent, data-backed bounds wherever the available results allow it.

\subsection{Quantitative uncertainty bounds (derived from Module~1 results)}
\label{subsec:error_quant}

\subsubsection{Baseline model-to-SCADA mismatch}
A direct, empirical bound is obtained by comparing the annualized SCADA AEP with the wake-aware baseline produced by the Module~1 model. Let $\mathrm{AEP}_{\mathrm{SCADA}}$ be the annualized measured production and $\mathrm{AEP}_{\mathrm{M1,wake}}$ the modeled production including wakes. The relative AEP mismatch is
\begin{equation}
\varepsilon_{\mathrm{AEP}} \;=\;
\frac{\left|\mathrm{AEP}_{\mathrm{M1,wake}}-\mathrm{AEP}_{\mathrm{SCADA}}\right|}{\mathrm{AEP}_{\mathrm{SCADA}}}.
\label{eq:err_aep_mismatch}
\end{equation}
Using the project results ($\mathrm{AEP}_{\mathrm{M1,wake}}=89.04~\mathrm{GWh}$ and $\mathrm{AEP}_{\mathrm{SCADA}}=84.18~\mathrm{GWh}$), the absolute deviation is $4.86~\mathrm{GWh}$ and
\begin{equation}
\varepsilon_{\mathrm{AEP}} \approx 0.058 \;\;\;(\text{about }5.8\%).
\label{eq:err_aep_value}
\end{equation}
This mismatch aggregates several effects into one empirical bound: wake-model simplifications (Jensen/Park), discretization choices, power-curve approximation, and real operational factors present in SCADA (e.g., curtailment, availability, and measurement noise) that are not fully represented in the simplified physics model.

\subsubsection{Wake-loss estimate sensitivity (model vs. SCADA)}
Wake loss is computed relative to the no-wake potential from Module~1, $\mathrm{AEP}_{\mathrm{M1,pot}}$, as
\begin{equation}
L \;=\; 1 - \frac{\mathrm{AEP}_{\mathrm{wake}}}{\mathrm{AEP}_{\mathrm{pot}}}.
\label{eq:err_wakeloss_def}
\end{equation}
Applying \eqref{eq:err_wakeloss_def} to both the modeled wake AEP and the SCADA AEP gives:
\begin{align}
L_{\mathrm{M1}} &= 1 - \frac{\mathrm{AEP}_{\mathrm{M1,wake}}}{\mathrm{AEP}_{\mathrm{M1,pot}}}
= 1 - \frac{89.04}{97.49} = 0.0867 \;\; (8.67\%), \label{eq:err_wakeloss_model}\\
L_{\mathrm{SCADA}} &= 1 - \frac{\mathrm{AEP}_{\mathrm{SCADA}}}{\mathrm{AEP}_{\mathrm{M1,pot}}}
= 1 - \frac{84.18}{97.49} = 0.1365 \;\; (13.65\%). \label{eq:err_wakeloss_scada}
\end{align}
The difference,
\begin{equation}
\Delta L_{\mathrm{pp}} \;=\; 100\left(L_{\mathrm{SCADA}}-L_{\mathrm{M1}}\right) \approx 5.0 \;\text{pp},
\label{eq:err_wakeloss_gap}
\end{equation}
indicates that the simplified wake model likely underestimates wake losses under real operation, or that the assumed no-wake potential overestimates the truly achievable reference. For this reason, the optimization outcomes are evaluated against the validated wake-aware baseline (Module~1) and interpreted primarily as \emph{wake-loss recovery within the model assumptions}.

\subsubsection{Wind-rose aggregation vs. time-series evaluation (Module~2 calibration residual)}
Module~2 uses a wind-rose aggregation, while Module~1 replays time-series SCADA. To make the scales comparable, Module~2 introduces the calibration factor $\eta$ (cf. \eqref{eq:m2_calibration}). A remaining consistency check is the post-calibration baseline residual:
\begin{equation}
\varepsilon_{\mathrm{base}} \;=\;
\left|
\frac{\mathrm{AEP}(\mathbf{x}_0)-\mathrm{AEP}_{\mathrm{M1,wake}}}{\mathrm{AEP}_{\mathrm{M1,wake}}}
\right|,
\label{eq:err_baseline_residual}
\end{equation}
which is targeted below $0.5\%$ in this project. This does not remove all structural errors (e.g., wake shape, terrain effects), but it avoids reporting artificial gains caused purely by different AEP evaluation conventions.

\subsection{Qualitative uncertainty sources (structural and data-driven)}
\label{subsec:error_qual}

Beyond the quantitative bounds above, several qualitative factors can bias absolute AEP and wake-loss estimates:

\begin{itemize}
    \item \textbf{Physical simplifications (2D, terrain, stability):} The farm is modeled in a 2D horizontal plane with simplified inflow. Real farms experience complex terrain speed-up/slow-down, wind veer, wake meandering, and atmospheric stability effects, which can change wake recovery and effective turbulence levels \cite{shakoor2016jensenreview}.
    \item \textbf{Wake model form (Jensen/Park limitations):} Jensen/Park assumes a linearly expanding, top-hat wake deficit and typically uses constant parameters ($k$, $C_T$). Real wakes are closer to Gaussian profiles and depend on turbulence intensity, stability, and shear; thus the wake deficit magnitude and overlap physics can deviate systematically from Jensen predictions \cite{katic1987jensen,shakoor2016jensenreview}.
    \item \textbf{Discretization effects:} (i) SCADA is evaluated at 10-minute averages, which can smooth nonlinear wake interactions; (ii) wind-rose sectoring approximates a continuous distribution, and wide directional bins can blur sharp wake-alignment events. This mainly affects absolute AEP, but relative comparisons remain meaningful if the same discretization is applied consistently.
    \item \textbf{SCADA data and operational states:} Curtailment, grid constraints, downtime, and yaw misalignment can reduce measured energy without being explicitly represented in the physics model. Filtering or imperfect classification of these states can shift the SCADA baseline and therefore the apparent model mismatch.
    \item \textbf{GIS data quality and buffer interpretation:} OSM completeness and geometric accuracy vary locally. Buffer distances are applied as hard constraints in a planar coordinate system; errors in feature geometry or CRS transformations can affect feasibility masks and the final siting envelope.
    \item \textbf{Stochastic optimization variability:} Differential Evolution and the sequential expansion step include stochastic sampling. With limited population size and iterations, different random seeds can yield slightly different local optima. A standard robustness check is to repeat optimization runs with multiple seeds and report the mean and standard deviation of achieved AEP under identical constraints \cite{herbertacero2014review}.
\end{itemize}

Overall, the key quantitative takeaway is that the baseline model reproduces annual SCADA AEP within about $5.8\%$ and that wake-loss estimates can differ by about $5$ percentage points when comparing SCADA-derived versus Jensen-based wake loss under the same no-wake reference. Consequently, the strongest claim supported by this work is the \emph{relative improvement} of optimized layouts compared to the validated wake-aware baseline under consistent assumptions.


\section{National-Scale Scalability Demonstration: Project Deutschland 2030}
\label{sec:germany2030}

To assess scalability and external validity, we formulated a coarse national expansion problem for Germany. Wind power contributed roughly one third of public electricity generation in 2024. According to preliminary Bundesnetzagentur figures, installed onshore wind capacity totalled $63.5~\mathrm{GW}$ at the end of 2024, with a policy target of $115~\mathrm{GW}$ by 2030 \cite{bnetza2025renewables2024}. This section evaluates whether the proposed optimization logic can be transferred from a single wind farm to a national-scale setting using clustering and proxy models.

\subsubsection{Data basis and clustering}
We used the German Marktstammdatenregister (MaStR) export of operational onshore wind units (coordinates and rated power) \cite{bnetza_mastr}. Since optimizing tens of thousands of individual turbines is computationally prohibitive, the fleet was aggregated via DBSCAN clustering (radius $2.5~\mathrm{km}$, minimum 3 turbines). Each cluster $i$ is represented by a centroid position $\mathbf{x}_i$ and an aggregated capacity $C_i$ (MW). Isolated turbines were retained for visualization only.

\subsubsection{Proxy wind resource model}
A simplified wind-speed proxy was defined as a north--south gradient (in EPSG:3035 northing $y$), inspired by spatial wind resource patterns visible in continental wind atlases \cite{globalwindatlas}. The proxy speed is
\begin{equation}
v(y)=v_{\min}+\left(\frac{y-y_{\min}}{y_{\max}-y_{\min}}\right)\left(v_{\max}-v_{\min}\right),
\end{equation}
with $(v_{\min},v_{\max})=(5.5,9.0)\,\mathrm{m/s}$ and $(y_{\min},y_{\max})=(2.6,3.6)\times 10^6~\mathrm{m}$.
The gross yield factor was approximated by the cubic dependency
\begin{equation}
\eta_i = \left(\frac{v(y_i)}{v_{\mathrm{ref}}}\right)^3,\qquad v_{\mathrm{ref}}=7.5~\mathrm{m/s}.
\end{equation}

\subsubsection{Cluster-level wake-loss proxy}
Wake interaction between clusters was modeled using a Jensen-type geometric deficit formulation \cite{shakoor2016jensenreview,katic1987jensen}. For each wind direction $d$ with probability $p_d$, cluster coordinates are rotated into the flow frame and wake influence is applied only downstream within a wake cone of effective diameter $D_c$.
The pairwise deficit contribution is
\begin{equation}
\delta_{j\leftarrow i}(d)=\left(1-\sqrt{1-C_T}\right)\left(\frac{D_c}{D_c+2k\Delta x_{ij}}\right)^2,
\end{equation}
if $\Delta x_{ij}>0$ and $|\Delta y_{ij}|<D_c/2+k\Delta x_{ij}$. Deficits are combined in quadratic sum. A capacity-weighted wake-loss fraction $L\in[0,1]$ is computed over all clusters and directions.

\subsubsection{Optimization objective, constraints, and expansion loop}
The objective was to maximize a net yield proxy:
\begin{equation}
Y = \left(\sum_i C_i \eta_i\right)(1-L),
\end{equation}
which is optimized by Differential Evolution over the location of one new cluster at a time. Each new site adds a fixed nominal capacity of $600~\mathrm{MW}$. Candidate sites were constrained to (i) lie inside Germany, (ii) satisfy a minimum inter-cluster distance of $15~\mathrm{km}$, and (iii) avoid simplified exclusion zones (buffers around major cities, lakes, and mountainous regions) to mimic land-use constraints.

The sequential expansion terminates once the nominal capacity target of $115~\mathrm{GW}$ is reached. Fig.~\ref{fig:germany2030} shows the resulting new cluster locations and the estimated wind-speed proxy.

\begin{figure}[H]
\centering
\includegraphics[width=\columnwidth]{germany_final_corrected_aep.png}
\caption{Germany 2030 scalability test: existing turbine clusters (purple), single turbines (cyan), and proposed new cluster sites (stars) colored by estimated wind-speed proxy. The sequential expansion reached the $115~\mathrm{GW}$ target (achieved $115.1~\mathrm{GW}$).}
\label{fig:germany2030}
\end{figure}

% ============================
% Conclusion
% ============================
\section{Conclusion and Outlook}
\label{sec:conclusion}

This paper presented a proof-of-concept workflow to improve wind farm performance by explicitly considering wake effects while respecting real land-use constraints. The work was structured into three modules: Module~1 established a baseline using 2023 SCADA data and quantified wake-related losses for the Hill of Towie wind farm; Module~2 optimized turbine positions and explored constrained expansion; and Module~3 verified that the reported improvements are comparable to the SCADA-based baseline and remain feasible under the defined GIS and spacing rules.

Overall, the results indicate that meaningful energy gains are achievable within a consistent modeling framework. In Module~2, the constrained optimization improved the calibrated AEP relative to the as-built baseline, and the dynamic expansion strategy demonstrated that additional turbines can be placed in a controlled way using explicit stopping criteria to avoid inefficient over-densification. The accompanying GIS plots confirm that the final layouts respect the defined boundary and exclusion zones.

As an outlook beyond a single wind farm, the national-scale scalability demonstration (\textit{Project Deutschland 2030}) suggests that the same general logic can be transferred to much larger settings. By clustering the German onshore fleet and using simplified proxy models, the sequential expansion reached the nominal 2030 capacity target, illustrating how wake-aware, constraint-based placement can be formulated at system level. While this national experiment is intentionally coarse, it serves as a practical indication that the approach can scale and can support larger planning questions.

The main limitation of this study is the simplified physics and data representation used for computational efficiency. Future work should improve realism (e.g., better wake and terrain modeling, refined constraint sets, and repeated robustness runs) and extend the objective from energy-only metrics toward planning-relevant criteria such as costs, grid connection, and acceptance-related constraints.



% pull a few lines from the last page back
\enlargethispage{3.5\baselineskip}

\bibliographystyle{IEEEtran}
\bibliography{references}
\end{document}
